% swim-times-srs.tex
% Software Requirements Specification for the swim-times program
% IEEE/ISO/IEC/IEEE 29148-2018 structure
% Compiled with: pdflatex swim-times-srs.tex

\documentclass[12pt]{article}

\usepackage{geometry}
\geometry{letterpaper, margin=1in}

\usepackage{hyperref}
\usepackage{booktabs}
\usepackage{array}
\usepackage{longtable}
\usepackage{enumitem}
\usepackage{titlesec}
\usepackage{fancyhdr}
\usepackage{xcolor}

\definecolor{castblue}{RGB}{0,48,135}

\pagestyle{fancy}
\fancyhf{}
\renewcommand{\headrulewidth}{0.4pt}
\fancyhead[L]{\small College Area Swim Team}
\fancyhead[R]{\small SRS: swim-times}
\fancyfoot[C]{\thepage}

\titleformat{\section}{\large\bfseries\color{castblue}}{\thesection}{0.5em}{}[\titlerule]
\titleformat{\subsection}{\normalsize\bfseries}{\thesubsection}{0.5em}{}

\newcommand{\code}[1]{\texttt{#1}}
\newcommand{\req}[1]{\textbf{[#1]}}

\begin{document}

% -----------------------------------------------------------------
% Title block
% -----------------------------------------------------------------
\begin{titlepage}
  \centering
  \vspace*{2cm}
  {\Huge\bfseries\color{castblue} Software Requirements Specification\\[0.5em]}
  {\Large swim-times: Automated Swimmer-Times Retrieval System}\\[2em]
  {\large College Area Swim Team (CAST)}\\[0.5em]
  {\normalsize Prepared for: CAST Coaching Staff and Meet Directors}\\[0.5em]
  {\normalsize Prepared by: CAST Technology Working Group}\\[2em]
  \begin{tabular}{ll}
    Document identifier: & CAST-SRS-001                                     \\
    Revision:            & 6.0                                              \\
    Date:                & September 18, 2026                              \\
    Status:              & Released                                         \\
    Governing standard:  & ISO/IEC/IEEE 29148-2018 (Requirements Engineering) \\
  \end{tabular}
  \vfill
  {\small\textit{This document is prepared in accordance with the principles
  of rational documentation articulated in D.\,L.\,Parnas and P.\,C.\,Clements,
  ``A Rational Design Process: How and Why to Fake It,''
  \textit{IEEE Transactions on Software Engineering}, vol.\,SE-12, no.\,2,
  pp.\,251--257, February 1986.}}
\end{titlepage}

\tableofcontents
\clearpage

% =================================================================
\section{Introduction}
% =================================================================

\subsection{Purpose}

This Software Requirements Specification (SRS) defines the
functional and non-functional requirements for the
\textbf{swim-times} program.  It is the binding agreement between
the requesting stakeholders (CAST coaching staff) and the
implementing developers, and it is the authoritative reference
against which the software shall be verified.

\subsection{Scope}

The software product specified by this document is named
\textbf{swim-times}.  It is a single-binary command-line utility
that:

\begin{itemize}
  \item Accepts swimmer- and event-selection options on the command line.
  \item Resolves each selected swimmer's USA~Swimming \code{PersonKey}
        via the public Sisense JAQL person-search endpoint.
  \item Retrieves all recorded swim times for each selected
        swimmer--event pair from the USA~Swimming Times Elasticube.
  \item Sorts results fastest-first and prints them as either a
        formatted table or comma-separated values.
\end{itemize}

The program does \textbf{not} write to any database, send email,
post to any social channel, or modify any USA~Swimming record.

\subsection{Definitions, acronyms, and abbreviations}

\begin{description}[style=nextline,leftmargin=2.5cm,labelindent=0pt]
  \item[CAST] College Area Swim Team.
  \item[ConOps] Concept of Operations document
                (CAST-CONOPS-001).
  \item[CSV] Comma-Separated Values.
  \item[CWEB] Knuth's literate-programming toolchain for C
              (\code{ctangle} + \code{cweave}).
  \item[ElastiCube] Sisense's columnar in-memory analytics database.
  \item[JAQL] JSON Analytics Query Language.
  \item[LCM] Long-Course Meters.
  \item[PersonKey] USA~Swimming's internal numeric swimmer identifier.
  \item[SCY] Short-Course Yards.
  \item[SDI~AGC] San Diego-Imperial Age-Group Championship.
  \item[SRS] Software Requirements Specification (this document).
\end{description}

\subsection{References}

\begin{enumerate}[label={[\arabic*]}]
  \item ISO/IEC/IEEE 29148-2018, \textit{Systems and software engineering---
        Life cycle processes---Requirements engineering}.
  \item D.\,L.\,Parnas and P.\,C.\,Clements, ``A Rational Design Process:
        How and Why to Fake It,'' \textit{IEEE Trans.\ Software Eng.},
        vol.\,SE-12, no.\,2, pp.\,251--257, Feb.\ 1986.
  \item CAST-CONOPS-001 Rev.\,1.0, \textit{Concept of Operations:
        swim-times}, April 21, 2026.
  \item \code{swim2/swim-times.w} --- CWEB literate program (source of record).
  \item D.\,E.\,Knuth, ``Literate Programming,''
        \textit{The Computer Journal}, vol.\,27, no.\,2,
        pp.\,97--111, 1984.
\end{enumerate}

\subsection{Overview}

Section~2 gives an overall description of the product, its users,
constraints, and assumptions.  Section~3 enumerates the specific
requirements, each tagged with a unique identifier of the form
\code{REQ-X-nn} so that requirements can be traced individually
through verification.  Section~4 specifies how each requirement
shall be verified.

% =================================================================
\section{Overall Description}
% =================================================================

\subsection{Product perspective}

swim-times is a self-contained, end-user command-line tool.
It depends on:

\begin{itemize}
  \item The USA~Swimming Sisense analytics endpoint at
        \texttt{usaswimming.sisense.com} for all data.
  \item \code{libcurl} for HTTPS transport.
  \item A POSIX C runtime for standard library functions.
\end{itemize}

It does not interoperate with any local database, message queue,
or other process.  It is invoked as a child of an interactive
shell or a build script.

\subsection{Product functions}

At a high level, swim-times performs:

\begin{enumerate}
  \item Command-line option parsing (\code{-o}, \code{-e}).
  \item Person-search HTTPS POST to resolve a swimmer's
        \code{PersonKey} from a name substring.
  \item Times-query HTTPS POST per (swimmer, event) pair.
  \item In-memory insertion sort by Sisense \code{SortKey}.
  \item Formatted table or CSV output to standard output.
  \item Diagnostic reporting to standard error.
\end{enumerate}

\subsection{User characteristics}

The program has two user classes:

\begin{description}[style=nextline]
  \item[Coaches and parents.]
        Comfortable with a terminal but not necessarily with
        programming.  They use only the documented command-line
        flags and do not modify the source.
  \item[Maintainers.]
        Familiar with C, CWEB, and \code{libcurl}.  They retangle
        and recompile when the Sisense bearer token rotates or
        when new swimmers, events, or output modes are added.
\end{description}

\subsection{Constraints}

\begin{description}[style=nextline]
  \item[Network.]
        The host running swim-times must have outbound HTTPS access
        to \texttt{usaswimming.sisense.com}.
  \item[Authentication.]
        A valid Sisense bearer token must be embedded in the
        compiled binary.  Token rotation requires a recompile.
  \item[Data source exclusivity.]
        Only the USA~Swimming data hub is consulted.
        Swimcloud is explicitly excluded by stakeholder direction.
  \item[Single-binary distribution.]
        The product is one statically describable executable;
        no installer, no configuration files.
  \item[Literate source of record.]
        The CWEB \code{.w} file is canonical; the generated
        \code{.c} file is a build artefact.
\end{description}

\subsection{Assumptions and dependencies}

\begin{itemize}
  \item \textbf{Live retrieval requires a signed-in caller.}  As of
        2026 the USA~Swimming data hub answers reference queries
        anonymously but refuses \code{GetMembersForFilters},
        \code{BestTimes}, and every other search or times endpoint with
        \code{403 Forbidden} unless \code{Usas-Sub-Id} names a
        signed-in subject.  Revision 1.0 of this document assumed
        anonymous access; that assumption is withdrawn.  See
        CAST-DR-001.
  \item \textbf{A USA~Swimming account is available.}  Revision 4.0
        makes signing in the program's default and only online
        behaviour (\req{REQ-N-01}); the operator is assumed to hold an
        account whose credentials are stored per \req{REQ-N-02}.
  \item \textbf{A newly-issued session must be activated} before the
        times API will honour it (\req{REQ-N-05}).  This is an observed
        property of the service, not a documented one, and may change
        without notice; \req{REQ-G-01} is what would detect the
        change.
  \item \code{GET .../TimesSearch/GetMember/$\langle$memberId$\rangle$}
        and the \code{SearchFilter/*} reference feeds remain open to
        anonymous callers.  This is not guaranteed to persist and the
        program shall detect its loss rather than assume it
        (\req{REQ-G-01}).
  \item Each \code{BestTimes} record continues to carry the fields
        \code{eventCode}, \code{swimTime}, \code{swimDate},
        \code{timeStandard}, and \code{meetName}.
  \item The local Postgres mirror is the system of record for any
        deployment without USA~Swimming credentials; \code{-o offline}
        is its read path.
  \item No swimmer of interest has more than 200 recorded times in
        any single event (matches the \code{MAX\_TIMES} compile-time
        constant).
\end{itemize}

% =================================================================
\section{Specific Requirements}
% =================================================================

\subsection{Functional requirements}

\subsubsection{Command-line interface}

\begin{description}[style=nextline]
  \item[\req{REQ-F-01}] The program shall accept the option flag
        \code{-o} followed by a comma-separated list of one or more
        tokens.  Each token shall be either a \emph{behaviour keyword}
        (\code{fastest}, \code{csv}, \code{store}, \code{offline}) or a
        \emph{swimmer id} (alias) drawn from the roster of
        \S\ref{sec:roster}.  Behaviour keywords and swimmer ids may be
        freely intermixed within a single \code{-o} list.
  \item[\req{REQ-F-02}] The program shall accept the option flag
        \code{-e} followed by one or more comma-separated
        USA~Swimming event codes, and shall accept the \code{-e}
        flag repeated multiple times.
  \item[\req{REQ-F-03}] When invoked with no options the program
        shall print a usage message to standard error and exit
        with status~1.
  \item[\req{REQ-F-04}] When invoked with an unrecognised flag the
        program shall print the usage message and exit with status~2.
\end{description}

\subsubsection{Swimmer selection}

\begin{description}[style=nextline]
  \item[\req{REQ-F-10}] Each swimmer in \texttt{SWIMMERS[]} shall
        carry a unique string \emph{id} (alias).  When one or more
        swimmer ids are supplied to \code{-o}, the program shall
        process only the swimmers whose id matches one of those
        supplied.
  \item[\req{REQ-F-11}] When no swimmer id is supplied to \code{-o}
        (whether or not events are supplied to \code{-e}), the
        program shall process every entry in \texttt{SWIMMERS[]} ---
        the full roster of \S\ref{sec:roster}.
  \item[\req{REQ-F-12}] The set of selectable swimmers and their ids
        (aliases) shall be exactly the roster enumerated in
        \S\ref{sec:roster} (\req{REQ-F-17}).  An unrecognised
        \code{-o} token that matches no swimmer id and no behaviour
        keyword shall select no swimmer and shall not be an error.
  \item[\req{REQ-F-13}] When any of the \code{ledecky10},
        \code{ledecky12}, or \code{ledecky14} keywords is
        supplied, the program shall resolve Katie Ledecky's
        \code{memberId} via the same member-search path used for
        the other swimmers.
  \item[\req{REQ-F-14}] When the \code{ledecky10} keyword is
        supplied, the program shall record the inclusive ISO date
        window \code{2006-03-17} through \code{2008-03-16}
        (Katie Ledecky's ninth and tenth years) on her
        \code{Swimmer} entry.
  \item[\req{REQ-F-15}] When the \code{ledecky12} keyword is
        supplied, the program shall record the inclusive ISO date
        window \code{2008-03-17} through \code{2010-03-16}
        (Katie Ledecky's eleventh and twelfth years).
  \item[\req{REQ-F-16}] When the \code{ledecky14} keyword is
        supplied, the program shall record the inclusive ISO date
        window \code{2010-03-17} through \code{2012-03-16}
        (Katie Ledecky's thirteenth and fourteenth years).
  \item[\req{REQ-F-16a}] The age-window fields of
        \req{REQ-F-14}--\req{REQ-F-16} are applied to the
        \code{offline} \texttt{SELECT} only.  The online
        \code{BestTimes} query returns a single lifetime best per
        course, so an online \code{ledecky*} run reports that best
        (with its swim date) rather than an age-filtered subset.
        Age-bracket filtering of the online path is deferred pending
        a signed-in richer-times endpoint.
\end{description}

\subsubsection{Swimmer roster (names and aliases)}
\label{sec:roster}

\begin{description}[style=nextline]
  \item[\req{REQ-F-05}] The usage text shall document every behaviour
        keyword (\code{fastest}, \code{csv}, \code{store},
        \code{offline}, \code{diag}, \code{selftest}), the \code{-m}
        flag, and the environment variables \code{USAS\_SUB\_ID} and
        \code{USAS\_SESSION\_ID}.
  \item[\req{REQ-F-06}] The roster listing in the usage text shall be
        generated from the \code{SWIMMERS} table, so that it cannot
        drift out of step with the roster it describes.
  \item[\req{REQ-F-07}] \code{diag} and \code{offline} shall be
        rejected as mutually exclusive, with exit status
        \code{RC\_USAGE}.
  \item[\req{REQ-F-18}] The program shall accept
        \code{-m $\langle$memberId$\rangle$}, which fetches that member
        directly through \code{GetMember} and bypasses both the roster
        and the member search.  A roster entry carrying a known
        \code{member\_id} shall take the same short cut automatically.
  \item[\req{REQ-F-17}] The program shall be designed to fetch times
        for exactly the swimmers listed in Table~\ref{tab:roster}.
        Each row gives the swimmer's full name, the unique id (alias)
        by which she is selected on the command line, and the
        person-search query and lower-cased match substring used to
        resolve her USA~Swimming \code{PersonKey}.  ``Evans, Stella~J''
        in the source roster request is the same swimmer as
        \code{stella} and is not duplicated.
\end{description}

\begin{longtable}{@{}l l l l@{}}
  \caption{Swimmer roster: names and aliases.}\label{tab:roster}\\
  \toprule
  \textbf{Full name} & \textbf{id (alias)} & \textbf{Search query}
    & \textbf{Match substr} \\
  \midrule
  \endfirsthead
  \toprule
  \textbf{Full name} & \textbf{id (alias)} & \textbf{Search query}
    & \textbf{Match substr} \\
  \midrule
  \endhead
  \bottomrule
  \endfoot
  Stella Julianna Evans   & \code{stella}              & Julianna Evans  & stella \\
  Kalea Rose Benavente    & \code{kalea}               & Benavente       & kalea \\
  Kenneth Ray Evans       & \code{kenny}               & Ray Evans       & kenneth ray \\
  Keith Santiago Evans    & \code{keith}               & Santiago Evans  & keith santiago \\
  Katie Genevieve Ledecky & \code{ledecky10}           & Ledecky         & katie \\
  Katie Genevieve Ledecky & \code{ledecky12}           & Ledecky         & katie \\
  Katie Genevieve Ledecky & \code{ledecky14}           & Ledecky         & katie \\
  Emma Katherine Bothwell & \code{bothwell-emma}       & Emma Bothwell        & emma \\
  Wally Baris Coombs      & \code{coombs-wally}        & Wally Coombs         & wally \\
  Eva Cruz Rae Cruz       & \code{cruz-eva}            & Eva Cruz             & cruz rae \\
  Zoe Elizabeth Darrow    & \code{darrow-zoe}          & Zoe Darrow           & zoe \\
  Eva Doan                & \code{doan-eva}            & Eva Doan             & eva doan \\
  Nova Larue Ellis        & \code{ellis-nova}          & Nova Ellis           & nova \\
  Liliya Penny Fedyshyn   & \code{fedyshyn-liliya}     & Liliya Fedyshyn      & liliya \\
  Layla Emmanuelle Glass  & \code{glass-layla}         & Layla Glass          & layla \\
  Logan Wiley Glass       & \code{glass-logan}         & Logan Glass          & wiley \\
  Hannah Francis Gosser   & \code{gosser-hannah}       & Hannah Gosser        & gosser \\
  Diamond Khaan Huynh     & \code{huynh-diamond}       & Diamond Huynh        & diamond \\
  Hannah Leigh Knuth      & \code{knuth-hannah}        & Hannah Knuth         & knuth \\
  Cleo L Lefebvre-Bailey  & \code{lefebvre-bailey-cleo}& Cleo Lefebvre-Bailey & cleo \\
  Gia N Mendez            & \code{mendez-giavanna}     & Giavanna Mendez      & mendez \\
  Sofia Elizabeth Mulvaney& \code{mulvaney-sofia}      & Sofia Mulvaney       & sofia \\
  Naia Lamb Shay          & \code{shay-naia}           & Naia Shay            & naia \\
  Briar Reece Simmons     & \code{simmons-briar}       & Briar Simmons        & briar \\
  Presley Aspen Simmons   & \code{simmons-presley}     & Presley Simmons      & aspen \\
  Lucas Erick Soto        & \code{soto-lucas}          & Lucas Soto           & erick \\
  Vivi Edith Soto         & \code{soto-vivienne}       & Vivienne Soto        & vivi \\
  Addison Marie Wittmershaus & \code{wittmershaus-addison}& Addison Wittmershaus & addison \\
  Elsie Marguerite Zahner & \code{zahner-elsie}        & Elsie Zahner         & elsie \\
\end{longtable}

\subsubsection{Event selection}

\begin{description}[style=nextline]
  \item[\req{REQ-F-20}] When one or more event codes are supplied
        to \code{-e}, the program shall query only those events.
  \item[\req{REQ-F-21}] When \code{-e} is omitted, the program
        shall query all thirty-one default events
        (fourteen SCY, seventeen LCM) listed in the swim-times ConOps.
  \item[\req{REQ-F-22}] The program shall support up to
        \code{NUM\_EVENTS} event codes per invocation; additional
        codes beyond that limit shall be silently dropped.
\end{description}

\subsubsection{Data retrieval}

\begin{description}[style=nextline]
  \item[\req{REQ-F-30}] For each selected swimmer the program
        shall resolve a \code{memberId} by issuing one HTTPS POST
        to \code{times-api.usaswimming.org/swims/TimesSearch/%
        GetMembersForFilters} (body \code{\{"name":"\dots"\}}) and
        matching the configured lower-cased substring against the
        returned \code{fullName} field.
  \item[\req{REQ-F-31}] For each distinct stroke-and-distance pair
        among the requested events the program shall issue one HTTPS
        POST to \code{.../swims/TimesSearch/BestTimes} (body
        \code{\{"memberId":"\dots","strokeAbbreviation":"\dots",%
        "distance":n\}}), which returns the swimmer's lifetime best
        for that pair in each course (SCY and~LCM).
  \item[\req{REQ-F-32}] The program shall read five fields per
        returned record: \code{eventCode}, \code{swimTime},
        \code{swimDate}, \code{timeStandard} (the motivational
        standard), and \code{meetName}.
  \item[\req{REQ-F-34}] A \code{404} from \code{BestTimes} shall be
        treated as ``this swimmer has no time in that stroke and
        distance'', not as an error: the event shall report
        \code{(no times found)} and the run shall continue.  For a
        young swimmer most of the thirty-one events answer this way.
  \item[\req{REQ-F-33}] The program shall cache all records returned
        for a member and answer every per-event query from that cache,
        issuing no more than one \code{BestTimes} POST per distinct
        stroke-and-distance pair.
\end{description}

\subsubsection{Result processing}

\begin{description}[style=nextline]
  \item[\req{REQ-F-40}] The program shall sort the rows for each
        event by swim time ascending (smallest = fastest), using the
        seconds value computed by \code{time\_to\_seconds}, before
        printing.
  \item[\req{REQ-F-41}] When the \code{fastest} keyword is supplied,
        the program shall print only the single fastest row per
        event, suppressing all slower entries.
  \item[\req{REQ-F-42}] The program shall reformat dates from the
        API's \code{"Mon DD, YYYY"} form to ISO \code{YYYY-MM-DD}.
\end{description}

\subsubsection{Output}

\begin{description}[style=nextline]
  \item[\req{REQ-F-50}] In the default (table) mode the program
        shall print, for each event, a heading line, a column
        header, a separator, and one row per recorded time
        containing time, date, motivational standard, and meet name.
  \item[\req{REQ-F-51}] When the \code{csv} keyword is supplied
        the program shall print exactly one CSV header line at the
        start of output and one CSV data line per recorded time.
        Each CSV data line shall contain swimmer name, event code,
        time, date, motivational standard, and meet name in that order.
  \item[\req{REQ-F-52}] When an event has zero recorded times in
        table mode the program shall print the literal
        \code{(no times found)} for that event.
\end{description}

\subsubsection{Error handling}

\begin{description}[style=nextline]
  \item[\req{REQ-F-60}] If a swimmer's \code{memberId} cannot be
        resolved the program shall print a diagnostic to standard
        error, free all libcurl global state, and exit with
        status~1.
  \item[\req{REQ-F-61}] If an individual times query fails, the
        program shall print a diagnostic to standard error and
        continue with the next event without aborting.
\end{description}

\subsubsection{Database persistence}

\begin{description}[style=nextline]
  \item[\req{REQ-F-70}] When the \code{store} keyword is supplied,
        the program shall write every emitted row to the local
        Postgres database \texttt{swim-times} programmatically via
        the \texttt{libpq} client library, using a prepared
        \texttt{INSERT \dots\ ON~CONFLICT~DO~NOTHING} per row.
        \code{store} composes with \code{csv}: stdout still
        receives the CSV stream while the rows are persisted.
  \item[\req{REQ-F-77}] A swimmer the USA~Swimming directory does not
        hold shall be skipped, not treated as a failure.  The member
        search answers \code{404} when no member matches the name, and
        the by-id lookup answers \code{404} for an unknown id; both mean
        ``there is no such swimmer'', which is a fact about one roster
        entry.  The run shall warn, count the entry, and continue with
        the next swimmer.  Any other non-\code{200} remains fatal.
  \item[\req{REQ-F-77a}] The count of skipped swimmers shall be
        repeated at the end of the run, where it cannot be lost in the
        output.  A run in which no swimmer at all resolved shall exit
        \code{RC\_UNAVAIL}; a run that fetched some and skipped others
        shall exit \code{RC\_OK}.
  \item[\req{REQ-F-78}] Before inserting a swimmer the program shall
        look for an existing row in three tiers: by \code{member\_id};
        by \code{full\_name} compared case-insensitively; and by first
        and last name only, likewise case-insensitively, which bridges
        the difference between a recorded middle initial and a returned
        middle name.  The second and third tiers shall apply
        \emph{only} to rows whose \code{member\_id} is \code{NULL}.  A
        row already carrying a different member id denotes a different
        swimmer however alike the names, and merging two swimmers is a
        worse outcome than duplicating one.
  \item[\req{REQ-F-79}] A loose match shall link the existing row
        rather than rename it: \code{member\_id}, \code{lsc\_code}, and
        \code{club\_name} are filled in where absent, and
        \code{full\_name} is left as the database already had it.
  \item[\req{REQ-F-80}] An insert refused because another swimmer
        already holds that name shall be reported, naming the swimmer,
        and her times skipped.  The program shall not attach one
        swimmer's times to another.
  \item[\req{REQ-F-76}] The offline read shall report every swim
        recorded for the swimmer in the requested event, joining
        \texttt{swim} to \texttt{swimmer} and, by outer join,
        to \texttt{meet}.  This includes rows written by the sibling
        loader and not only those this program wrote.
  \item[\req{REQ-F-71}] When the \code{offline} keyword is
        supplied, the program shall read swim rows from the local
        Postgres \texttt{swim-times} database and emit them in the
        same table or CSV form as the online path, without
        touching the network.
  \item[\req{REQ-F-72}] Under \code{offline}, the program shall
        not initialise libcurl and shall not issue any HTTP
        request.
  \item[\req{REQ-F-73}] The \texttt{swim\_times} schema shall
        enforce a UNIQUE constraint on
        \code{(swimmer, event, swim\_time, swim\_date, meet)} so
        that re-running a \code{store} invocation does not
        introduce duplicate rows.
  \item[\req{REQ-F-74}] If the libpq connection cannot be
        established for \code{store}, the program shall print a
        single warning to standard error and continue with the
        run, without persisting any rows.  Under \code{offline} a
        connection failure is fatal and exits with status~1.
  \item[\req{REQ-F-75}] When \code{store} is rerun against a
        \texttt{swim\_times} table that already contains the rows
        being inserted, the program shall recover gracefully: each
        duplicate row shall be silently absorbed by the
        \texttt{ON~CONFLICT~DO~NOTHING} clause without aborting
        the transaction, the process shall exit~0, and no
        diagnostic shall be written to standard error.
        (This realises requirement~\#4 of \texttt{requirements4.txt}.)
\end{description}

\subsection{Authentication requirements}

Added in revision 4.0.  They replace the credential-passing scheme of
revision 3.0, in which the operator was expected to extract a session
from a browser by hand.

\begin{description}[style=nextline]
  \item[\req{REQ-N-01}] Every run that is not \code{offline} shall sign
        in to the USA~Swimming data hub before issuing any query.  The
        program shall not attempt an unauthenticated search or times
        query, and shall not fall back to one when a signed-in request
        fails.
  \item[\req{REQ-N-02}] Credentials shall be read from the file named
        by \code{USAS\_ENV\_FILE}, defaulting to
        \code{\$HOME/.usas-env}, in shell assignment syntax
        (\code{[export] NAME="value"}).  The file shall be
        \emph{parsed}, never executed, and only names prefixed
        \code{USAS\_} shall be adopted --- such a file may carry
        settings intended for other programs, and adopting them would
        silently change this program's behaviour.
  \item[\req{REQ-N-03}] A value already present in the environment
        shall take precedence over the same name in the file.
  \item[\req{REQ-N-04}] The program shall obtain a subject and session
        id by completing the data hub's OpenID~Connect
        authorization-code flow: begin at the back-end-for-front-end,
        submit the identity provider's login form with the stored
        credentials and its antiforgery token, exchange the resulting
        authorization code, and read \code{sub} and \code{sid} from the
        user-info endpoint.
  \item[\req{REQ-N-05}] After obtaining a session the program shall
        call \code{POST security/auth/GetDataHubSecurityInfoForIdp}
        with that \code{sub} and \code{sid}.  Until this call is made
        the times API answers \code{401} for the new session; after it,
        \code{200}.
  \item[\req{REQ-N-06}] A sign-in failure shall be reported in terms of
        its cause --- absent credentials, rejected credentials, or a
        failed step --- and shall exit \code{RC\_NOPERM}.  In
        particular, credentials rejected by the identity provider shall
        not be reported as the \code{HTTP 200} the provider answers
        with.
  \item[\req{REQ-N-07}] \code{offline} shall require no credentials and
        shall attempt no sign-in.
  \item[\req{REQ-N-08}] No credential shall be written to standard
        output, to the database, or to any file.  The greeting naming
        the signed-in user shall be suppressed in \code{csv} mode,
        where it would corrupt the output.
\end{description}

\subsection{Diagnostic requirements}

\begin{description}[style=nextline]
  \item[\req{REQ-G-01}] The program shall provide a \code{-o diag}
        mode that probes at least one endpoint from each data-hub
        access class --- reference feed, single-member lookup, member
        search, and times query --- and reports each probe's HTTP
        status.  It shall exit \code{RC\_OK} regardless of what the
        probes return, since reporting a refusal is a successful
        diagnosis.
  \item[\req{REQ-G-02}] \code{-o diag} shall report the outcome of the
        sign-in attempt and the identity the program is presenting: the
        \code{Usas-Sub-Id} in force, whether a session id is set, and
        the \code{Device-Id} being sent.
  \item[\req{REQ-G-03}] The output of \code{-o diag} shall be
        sufficient, on its own, to distinguish an authorised client
        from absent credentials, from rejected credentials, from a
        signed-in caller the service still refuses, from an outage.
  \item[\req{REQ-G-04}] \code{-o diag} is the one mode that may proceed
        after a failed sign-in, since reporting what an unauthenticated
        caller can still reach is the diagnosis being asked for.  It
        shall label the probe as unauthenticated when it does so.
\end{description}

\subsection{Error-reporting and exit-status requirements}

\begin{description}[style=nextline]
  \item[\req{REQ-E-01}] Process exit status shall distinguish the
        categories of outcome: \code{0} success; \code{2}
        (\code{RC\_USAGE}) a command-line error; \code{69}
        (\code{RC\_UNAVAIL}) the data hub could not be reached or
        failed for a reason other than authorisation; \code{77}
        (\code{RC\_NOPERM}) the data hub refused the request for want
        of credentials.
  \item[\req{REQ-E-02}] Every diagnostic arising from an HTTP response
        shall name the operation attempted, the HTTP status with its
        reason phrase, and the URL.  No diagnostic shall assert a cause
        the program has not observed.
  \item[\req{REQ-E-03}] Remedial advice for an authorisation failure
        shall be printed at most once per run.
  \item[\req{REQ-E-04}] An authorisation failure shall abandon the run
        rather than reissue the same refused request for each remaining
        event and swimmer.
  \item[\req{REQ-E-05}] A transport failure shall be reported with
        libcurl's own description of the error, distinctly from an
        HTTP-level failure.
\end{description}

\subsection{Cryptographic requirements}

\begin{description}[style=nextline]
  \item[\req{REQ-S-01}] The \code{rate-key} header presented by a
        signed-in caller shall be computed as
        $\lfloor t/s \rfloor$, a period, and the lower-case
        hexadecimal HMAC-SHA256 of the decimal spelling of $t$ keyed by
        the session id, where $t$ is the Unix time in milliseconds
        divided by $10^{4}$ and $s$ is the number of seconds elapsed
        since midnight UTC.
  \item[\req{REQ-S-02}] SHA-256 and HMAC-SHA256 shall be implemented
        within the literate source.  No cryptographic library shall be
        introduced as a dependency (see \req{REQ-D-04}).
  \item[\req{REQ-S-03}] The program shall provide a \code{-o selftest}
        mode that verifies both primitives against the published
        vectors of FIPS~180-4 and RFC~4231 and exits non-zero if any
        vector fails.
\end{description}

\subsection{External interface requirements}

\subsubsection{User interfaces}

\begin{description}[style=nextline]
  \item[\req{REQ-I-01}] All user interaction shall be via
        command-line arguments and process exit status; the
        program shall not read from standard input.
\end{description}

\subsubsection{Software interfaces}

\begin{description}[style=nextline]
  \item[\req{REQ-I-10}] The program shall use the
        \code{libcurl} ``easy'' interface for all HTTP transport.
  \item[\req{REQ-I-11}] The program shall send the data-hub headers
        \code{Content-Type: application/json}, \code{AppName: DataHub},
        \code{Usas-Sub-Id}, and \code{Device-Id} on every request, plus
        \code{Origin} and \code{Referer} naming the data-hub front end.
        When a session id is configured --- which, after
        \req{REQ-N-01}, is every online request --- it shall
        additionally send \code{Usas-Session-Id} and \code{rate-key}.  It shall send no
        \code{Authorization} header (the service uses none) and no
        \code{Accept} header (see \req{REQ-H-05}).
  \item[\req{REQ-I-12}] The program shall reach Postgres
        \emph{programmatically} through the \texttt{libpq} client
        library --- never by spawning \texttt{psql} or any other
        external script (this realises requirement~\#3 of
        \texttt{requirements4.txt}).
  \item[\req{REQ-I-13}] The program shall use exactly one database,
        \texttt{swimming} on host \texttt{100.77.243.69}, named in a
        compiled-in constant.  There shall be no environment override:
        \texttt{SWIM\_TIMES\_PGCONNINFO} is withdrawn.  The database is
        shared with another loader, and a run that silently wrote
        elsewhere --- to a stale local copy, say --- would appear to have
        succeeded while contributing nothing.
  \item[\req{REQ-I-14}] The program shall adapt to that database's
        normalised schema rather than introduce a table of its own.
        Specifically: a swimmer shall be resolved to a
        \texttt{swimmer\_key} by \texttt{member\_id} \emph{or}
        \texttt{full\_name} (both are UNIQUE and some rows predate the
        member id), creating or refreshing the row as needed; a meet
        shall be resolved to a \texttt{meet\_key} by name, creating it
        if new; and a swim shall be written to \texttt{swim} with
        \texttt{source\_endpoint = 'BestTimes'} and
        \texttt{auth\_mode = 'authenticated'}.
  \item[\req{REQ-I-15}] Because \texttt{swim.swim\_time\_id} is a
        primary key with no default and \texttt{BestTimes} returns no
        \texttt{swimTimeId}, the program shall synthesise one: the
        leading 63 bits of the SHA-256 of the row's natural key,
        negated.  The value shall be stable across runs, so that a
        re-fetch reproduces it, and negative, so that it cannot collide
        with an id the service itself issued.
  \item[\req{REQ-I-16}] A motivational standard absent from the
        \texttt{time\_standard} reference table --- the elite labels
        \texttt{Nats}, \texttt{Trials}, \texttt{Summer Jrs} --- shall be
        written as SQL \texttt{NULL}, the foreign key admitting no
        other value.  The label shall still appear in the printed table
        and in the CSV, neither of which is constrained by the schema.
  \item[\req{REQ-I-17}] Writes shall be idempotent: a repeated fetch
        shall add no rows.  \texttt{ON CONFLICT DO NOTHING} without a
        target satisfies both the natural-key constraint and the
        primary key.
\end{description}

\subsubsection{Communication interfaces}

\begin{description}[style=nextline]
  \item[\req{REQ-I-20}] All network communication shall use HTTPS
        (TLS) to the host \texttt{times-api.usaswimming.org}.  No
        plaintext HTTP URL for any USA~Swimming host shall appear in
        the source.
\end{description}

\subsubsection{HTTP transport contract}

These requirements were added in revision 3.0 following CAST-DR-001.
Their absence is why the defect recorded there escaped detection: the
transport layer discarded the one observation that would have explained
the outage.

\begin{description}[style=nextline]
  \item[\req{REQ-H-01}] The HTTP request function shall return the
        response's HTTP status code to its caller.
  \item[\req{REQ-H-02}] That status shall be obtained from libcurl via
        \code{curl\_easy\_getinfo(\dots, CURLINFO\_RESPONSE\_CODE,
        \dots)} and not inferred from any other signal.
  \item[\req{REQ-H-03}] A response that completed but carried no body
        shall be returned as an empty string.  A null return shall mean
        that the request did not complete, and shall mean nothing else.
  \item[\req{REQ-H-04}] The request function shall set both a
        whole-transfer timeout and a connect timeout, so that an
        unresponsive host cannot stall a run indefinitely.
  \item[\req{REQ-H-05}] The program shall not send an \code{Accept}
        header.  The service's content negotiation responds to
        \code{Accept: application/json} by returning its payload
        re-encoded as a single escaped JSON string, which the response
        scanner cannot parse.
  \item[\req{REQ-H-06}] Every site that consumes an HTTP response shall
        check the status before parsing the body.
\end{description}

\subsection{Performance requirements}

\begin{description}[style=nextline]
  \item[\req{REQ-P-01}] A full refresh of all seven
        \texttt{SWIMMERS[]} entries across all thirty-one default
        events ($7 \times (1 + 31) = 224$ HTTP requests, including
        the three Ledecky age-window entries that share an
        underlying person but each issue their own person-search
        and event queries) shall complete within five minutes on a
        residential broadband connection.
  \item[\req{REQ-P-02}] A single-swimmer single-event query
        (\code{-o stella -e "100 FR SCY"}) shall complete within
        ten seconds under normal data-hub response time.
  \item[\req{REQ-P-03}] An offline single-swimmer single-event
        query (\code{-o stella,offline -e "100 FR SCY"}) shall
        complete within two seconds against a local Postgres
        instance carrying typical row counts.
\end{description}

\subsection{Design constraints}

\begin{description}[style=nextline]
  \item[\req{REQ-D-01}] The software shall be implemented in C
        conformant to ISO/IEC 9899:1999 (C99) or later.
  \item[\req{REQ-D-02}] The canonical source shall be a CWEB
        literate program (\code{swim-times.w}); the
        \code{.c} file shall be a build artefact only.
  \item[\req{REQ-D-03}] No source-code module (named CWEB chunk)
        shall exceed twenty-four lines, in keeping with the
        existing literate-programming style of the codebase.
  \item[\req{REQ-D-04}] No third-party dependency other than
        \code{libcurl}, \code{libpq} (the Postgres C client
        library, added by requirements4.txt~\#3), and the standard
        C library shall be linked into the binary.  Transitive
        TLS, compression, and authentication libraries pulled in
        by \code{libcurl} or \code{libpq} are also permitted.
  \item[\req{REQ-D-10}] The Postgres schema for persisted swim
        rows shall be defined in \texttt{schema.sql} adjacent to
        the source.  The table \texttt{swim\_times} shall expose
        the seven application columns listed in REQ-F-32 plus an
        identity primary key and a \texttt{fetched\_at} timestamp,
        and shall declare a UNIQUE constraint per REQ-F-73.
  \item[\req{REQ-D-11}] No credential --- bearer token, session id, or
        subject --- shall appear in the source text.  All credentials
        shall be read from the environment at run time.
\end{description}

\subsection{Software system attributes}

\begin{description}[style=nextline]
  \item[\req{REQ-A-01} (Reliability)] The program shall produce
        identical output for identical data-hub responses on
        repeated invocations.
  \item[\req{REQ-A-02} (Maintainability)] All data-hub endpoint URLs
        shall be declared as named constants in clearly identifiable
        CWEB chunks.  Credentials shall not be declared there at all
        (\req{REQ-D-11}).
  \item[\req{REQ-A-03} (Portability)] The program shall build
        and run unmodified on any POSIX.1-2008 system with
        \code{libcurl} 7.0 or later installed.
  \item[\req{REQ-A-04} (Self-documentation)] The CWEB source
        shall, when processed by \code{cweave} and \code{pdftex},
        produce a human-readable typeset document describing both
        the why and the how of the implementation.
\end{description}

% =================================================================
\section{Verification}
% =================================================================

Each requirement shall be verified by one of the following methods:

\begin{description}[style=nextline]
  \item[Inspection (I).]
        Visual examination of source, output, or documentation.
  \item[Demonstration (D).]
        Execution of the program with selected inputs and
        observation of the resulting output.
  \item[Test (T).]
        Execution of an automated test that compares actual
        output to an expected reference.
  \item[Analysis (A).]
        Reasoning about the program (e.g., line-count audit
        of CWEB chunks for \code{REQ-D-03}).
\end{description}

\medskip
\noindent
The verification matrix below assigns a primary method to each
requirement.  Test cases and demonstration scripts shall be
documented separately in the swim-times Verification Plan.

\begin{center}
\begin{longtable}{ll}
  \toprule
  \textbf{Requirement} & \textbf{Method} \\
  \midrule\endhead
  REQ-F-01 \ldots REQ-F-07 & Demonstration \\
  REQ-F-10 \ldots REQ-F-16 & Demonstration \\
  REQ-F-18                 & Demonstration \\
  REQ-F-20 \ldots REQ-F-22 & Demonstration \\
  REQ-F-30 \ldots REQ-F-33 & Inspection + Demonstration \\
  REQ-F-40 \ldots REQ-F-42 & Test \\
  REQ-F-50 \ldots REQ-F-52 & Demonstration \\
  REQ-F-60 \ldots REQ-F-61 & Demonstration \\
  REQ-F-70 \ldots REQ-F-75 & Test (DB-gated) \\
  REQ-I-01 \ldots REQ-I-20 & Inspection \\
  REQ-H-01 \ldots REQ-H-06 & Inspection of \texttt{swim-times.w} \\
  REQ-E-01 \ldots REQ-E-05 & Test (tier-gated) \\
  REQ-S-01 \ldots REQ-S-03 & Test (known-answer vectors) \\
  REQ-G-01 \ldots REQ-G-04 & Demonstration (\texttt{-o diag}) \\
  REQ-N-01, REQ-N-06, REQ-N-07 & Test (sign-in failure paths) \\
  REQ-N-02, REQ-N-03 & Test (fixture credentials file) \\
  REQ-N-04, REQ-N-05 & Demonstration (live sign-in) + Inspection \\
  REQ-N-08 & Inspection + Test (CSV cleanliness) \\
  REQ-F-34 & Demonstration (all-events run) \\
  REQ-I-13 \ldots REQ-I-17 & Inspection + Test (\texttt{test-db.sh}) \\
  REQ-F-76 & Test (joined offline read) \\
  REQ-F-77, REQ-F-77a & Test (absent swimmer) \\
  REQ-F-78 \ldots REQ-F-80 & Test (\texttt{test-db.sh} DB-19, DB-20) \\
  REQ-P-01, REQ-P-02, REQ-P-03 & Demonstration (timed) \\
  REQ-D-01, REQ-D-02, REQ-D-04 & Inspection \\
  REQ-D-03                 & Analysis (CWEB chunk line-count audit) \\
  REQ-D-10                 & Inspection of \texttt{schema.sql} \\
  REQ-D-11                 & Inspection (credential pattern scan) \\
  REQ-A-01                 & Test \\
  REQ-A-02, REQ-A-04       & Inspection \\
  REQ-A-03                 & Demonstration (build on second POSIX host) \\
  \bottomrule
\end{longtable}
\end{center}

\vfill
\begin{center}
  \small\textit{End of document CAST-SRS-001 Rev.\,6.0}
\end{center}

\end{document}
